\begin{table}[h!]
\centering
\caption{\textbf{Relationship Between Impacts and Device Experience (III to II)}}
\label{tab:impacts_vs_exp_3_2}
\begin{tabular}{lcc}
\hline
\hline
\textbf{Variable} & \textbf{(1)} & \textbf{(2)} \\
\midrule
Device Experience (years) & -0.0051* & -1.5231* \\
 & (0.0019) & (0.7125) \\
Intercept & 0.2379* & 82.2241+ \\
 & (0.1075) & (40.7092) \\
Baseline Serious Adverse Events & -0.0005 & -1.7597 \\
 & (0.0040) & (1.4963) \\
FDA's Own Accord & 0.0128 & -21.9656 \\
 & (0.0401) & (15.2212) \\
Orthopedic Devices & 0.0736 & 5.5511 \\
 & (0.0952) & (35.3210) \\
Other Devices & -0.1523+ & -52.1829+ \\
 & (0.0761) & (27.9667) \\
Plastic Surgery, Gastroenterology-Urology Devices & 0.0419 & -24.7850 \\
 & (0.0935) & (34.7247) \\
Baseline Patent Count & --- & 2.4315*** \\
 & --- & (0.5018) \\
\midrule
N & 25 & 25 \\
Serious-Event Estimate & Yes & No \\
Patenting Estimate & No & Yes \\
\bottomrule

\end{tabular}
\caption*{
\begin{footnotesize}Note: The regression relates device-level DiD estimates for (1) serious adverse events and (2) patenting activity to each device type's market age (FDA years of experience) and other covariates, including baseline patenting rates. \textquotedblleft---\textquotedblright\ indicates variables excluded from a given specification. Due to limited data, fixed effects are included only for higher-risk device categories (each with $\geq$ 2 observations); cardiovascular devices are the omitted group. To increase power, FDA-initiated reclassifications (own accord = 0) are included with a corresponding fixed effect. +, *, **, *** denote significance at 0.10, 0.05, 0.01, 0.001 levels.
\end{footnotesize}
}
\end{table}